\section{Posterior computation over ordered partitions}
\label{sec:exact-inference}

Once the prior-adjusted block array
$\widetilde A^{(r)}_{ij}=A^{(r)}_{ij}\gamma_u(i,j)$ has been constructed, the observation family no longer enters the partition recursions. This section gives the two distinct inference routes used by \BayesBreak: sum-product for posterior quantities and max-sum for a joint maximizing partition.

\begin{figure*}[t]
  \centering
  \resizebox{0.91\textwidth}{!}{
\begin{tikzpicture}[node distance=7mm and 12mm]
\node[secondarybox,text width=33mm] (score) {Same local log scores\\$\ell_{ij}$};
\node[primarybox,above right=of score,text width=36mm] (sp) {$\oplus=\logsumexp$\\sum-product};
\node[primarybox,below right=of score,text width=36mm] (ms) {$\oplus=\max$\\max-sum};
\node[successbox,right=of sp,text width=35mm] (marg) {Marginal posterior quantities};
\node[successbox,right=of ms,text width=35mm] (path) {One joint maximizing path};
\draw[arrPrimary] (score)--(sp);\draw[arrPrimary] (score)--(ms);
\draw[arrPrimary] (sp)--(marg);\draw[arrPrimary] (ms)--(path);
\draw[Warning,densely dashed,line width=.8pt] (marg.south)--node[right,font=\scriptsize,text=Warning]{not generally equal}(path.north);
\end{tikzpicture}
}
  \caption{The same local scores support two algebraically related but inferentially different routes. Sum-product integrates over paths; max-sum selects one globally maximizing path. Coordinatewise modes of marginal boundary distributions need not form that path.}
  \label{fig:semiring-routes}
\end{figure*}

\subsection{Forward and backward partition sums}
For $k\in\{0,\ldots,k_{\max}\}$ and $j\in\{0,\ldots,n\}$, let $L_{k,j}$ be the integrated score of all $k$-block partitions of the prefix $(0,j]$. Let $R_{k,i}$ be the corresponding score for all $k$-block partitions of the suffix $(i,n]$. Set
\begin{align*}
L_{0,0}&=1,\qquad L_{0,j}=0\quad(j>0),\\
R_{0,n}&=1,\qquad R_{0,i}=0\quad(i<n).
\end{align*}
The sum-product recursions are
\begin{align}
L_{k+1,j}&=\sum_{h=k}^{j-1}L_{k,h}\widetilde A^{(0)}_{h,j},
\label{eq:forward}\\
R_{k+1,i}&=\sum_{h=i+1}^{n-k}\widetilde A^{(0)}_{i,h}R_{k,h}.
\label{eq:backward}
\end{align}
In implementation, Equations~\eqref{eq:forward}--\eqref{eq:backward} are evaluated in log space using \texttt{logsumexp}.

\begin{theorem}[Partition-sum identities and boundary marginals]
\label{thm:sum-product}
Under Assumptions~\ref{ass:offline} and~\ref{ass:factorized-prior}, for every admissible $k$,
\begin{equation}
L_{k,n}=R_{k,0}
=\sum_{0=t_0<\cdots<t_k=n}\prod_{q=1}^{k}\widetilde A^{(0)}_{t_{q-1},t_q}.
\label{eq:partition-sum}
\end{equation}
Moreover, for $1\le p\le k-1$ and admissible $h$,
\begin{align}
p(t_p=h\mid y,k)
&=\frac{L_{p,h}R_{k-p,h}}{L_{k,n}},\nonumber\\
p(b_h=1\mid y,k)
&=\sum_{p=1}^{k-1}p(t_p=h\mid y,k).
\label{eq:boundary-marginals}
\end{align}
\end{theorem}
\begin{proof}
Induction on $k$ in Equation~\eqref{eq:forward} shows that $L_{k,j}$ sums the product of local scores over every ordered $k$-block path ending at $j$. Backward induction gives the same statement for $R_{k,i}$ on suffixes, proving Equation~\eqref{eq:partition-sum}. A partition satisfying $t_p=h$ has a unique decomposition into a $p$-block prefix ending at $h$ and a $(k-p)$-block suffix beginning at $h$. Their integrated scores multiply, and division by the total partition sum yields the first identity in Equation~\eqref{eq:boundary-marginals}. Strict ordering permits at most one boundary label $p$ at any fixed $h$ within a partition, so summing over $p$ gives the boundary-event probability.
\end{proof}

Let $p(k)$ be a segment-count prior. Because $C_k(u)$ normalizes Equation~\eqref{eq:partition-prior},
\begin{equation}
 p(k\mid y)\propto p(k)\frac{L_{k,n}}{C_k(u)},
 \qquad k=1,\ldots,k_{\max}.
\label{eq:k-posterior}
\end{equation}
The posterior may be reported in full; choosing $\widehat k=\arg\max_k p(k\mid y)$ is a separate decision rule. Conditional boundary distributions in Equation~\eqref{eq:boundary-marginals} and boundary-event probabilities can then be computed at fixed $k$ or averaged over $p(k\mid y)$.

\subsection{Posterior moments and Bayes curves}
For $r\in\{0,1,2\}$, define a block-covering contribution at fixed $k$,
\begin{equation}
F^{(r)}_{ij}(k)
=\frac{1}{L_{k,n}}
  \sum_{m=1}^{k}L_{m-1,i}\widetilde A^{(r)}_{ij}R_{k-m,j}.
\label{eq:block-cover}
\end{equation}
When the target at coordinate $s$ is the segment-level quantity $m_\star(\theta)$, its conditional posterior moments are
\begin{equation}
\mathbb E\!\left[m(\theta_s)^r\mid y,k\right]
 =\sum_{i<s\le j}F^{(r)}_{ij}(k).
\label{eq:bayes-curve}
\end{equation}
The posterior mean Bayes curve uses $r=1$; the $r=2$ result gives pointwise posterior variance after subtracting the squared mean.

\begin{proposition}[Block-covering decomposition]
\label{prop:block-cover}
If every admissible partition assigns each coordinate to exactly one block and $A^{(r)}_{ij}$ is the integrated $r$th moment numerator for that block, then Equation~\eqref{eq:bayes-curve} is the exact fixed-$k$ posterior moment of the segment-level target at $s$.
\end{proposition}
\begin{proof}
Every partition contains one and only one block $(i,j]$ satisfying $i<s\le j$. For a specified block position $m$, the score of all compatible prefixes is $L_{m-1,i}$, the integrated target contribution on $(i,j]$ is $\widetilde A^{(r)}_{ij}$, and the score of all compatible suffixes is $R_{k-m,j}$. Summing over $m$, then over all blocks covering $s$, enumerates each $k$-block partition exactly once with its integrated target numerator. Division by $L_{k,n}$ normalizes the posterior.
\end{proof}

\subsection{Joint MAP segmentation}
Define local log scores $\ell_{ij}=\log\widetilde A^{(0)}_{ij}$ and max-sum messages
\begin{equation}
M_{k,j}=\max_{h\in\{k-1,\ldots,j-1\}}
       \{M_{k-1,h}+\ell_{h,j}\},
\quad M_{0,0}=0,
\label{eq:max-sum}
\end{equation}
with $M_{0,j}=-\infty$ for $j>0$. Store a deterministic maximizing predecessor $h^\star(k,j)$. Starting at $(k,n)$ and following these predecessors recovers a complete boundary vector.

\begin{theorem}[Correctness of max-sum backtracking]
\label{thm:map-correctness-paper}
For fixed $k$, the boundary vector returned by Equation~\eqref{eq:max-sum} and its stored predecessors belongs to
\begin{equation}
\arg\max_{0=t_0<\cdots<t_k=n}
  \sum_{q=1}^{k}\ell_{t_{q-1},t_q}.
\label{eq:joint-map-objective}
\end{equation}
Consequently, it is a joint MAP segmentation conditional on $k$. Ties are resolved only by the declared predecessor convention.
\end{theorem}
\begin{proof}
For $k=1$, $M_{1,j}=\ell_{0,j}$ is the only one-block path score. Suppose $M_{k-1,h}$ is the optimal score among all $(k-1)$-block paths ending at $h$. Every $k$-block path ending at $j$ has a unique penultimate endpoint $h$ and total score equal to its prefix score plus $\ell_{h,j}$. Maximizing over $h$ therefore produces the optimal $k$-block score. Induction proves optimality at $(k,n)$. Each stored predecessor attains the corresponding maximum, so backtracking constructs a path attaining that terminal score.
\end{proof}

\begin{cautionbox}[Marginal modes are not a joint MAP path]
The vector formed from $\arg\max_h p(t_p=h\mid y,k)$ for each $p$ optimizes separate marginal decisions. Its entries can be incompatible or can have lower joint posterior probability than another partition. \BayesBreak therefore reports marginal modes as diagnostics and reserves ``joint MAP'' for Theorem~\ref{thm:map-correctness-paper}.
\end{cautionbox}

\subsection{Complexity and numerical invariants}
If cumulative sufficient statistics make each local conjugate score constant time, there are $\Theta(n^2)$ blocks to precompute. Sum-product and max-sum each inspect $\Theta(k_{\max}n^2)$ predecessor transitions. Values and backpointers require $\Theta(k_{\max}n)$ working storage, while retaining all block-cover terms can require $\Theta(k_{\max}n^2)$ storage.

\begin{proposition}[Complexity]
\label{prop:complexity-paper}
When cumulative sufficient statistics make each conjugate segment calculation constant time, segment precomputation costs $\Theta(n^2)$ time and storage; either partition recursion costs $\Theta(k_{\max}n^2)$ time. For a nonconjugate block routine with per-block cost $T_{\rm block}$, local precomputation costs $\Theta(n^2T_{\rm block})$ and the global recursions are unchanged.
\end{proposition}
\begin{proof}
There are $n(n+1)/2=\Theta(n^2)$ admissible intervals. For each $k$ and endpoint $j$, a recursion scans at most $j$ predecessors, so $\sum_{k=1}^{k_{\max}}\sum_{j=k}^{n}\Theta(j)=\Theta(k_{\max}n^2)$ when all blocks are admissible. Replacing constant-time local evaluation by cost $T_{\rm block}$ multiplies only the block-construction term.
\end{proof}

Four numerical identities are checked in every implementation: the normalized segment-count posterior sums to one; forward and backward total evidence agree; conditional boundary-event probabilities sum to $k-1$; and the backtracked path score equals the terminal max-sum score within numerical tolerance. These checks test algebraic implementation rather than empirical model quality.
